%==============================================================================
% Sjabloon projectvoorstel graduaatsproject
%==============================================================================
% Je gegevens staan in ../projectgegevens.tex. De inhoud van je voorstel
% schrijf je in voorstel-inhoud.tex - dat bestand wordt later ook automatisch
% als bijlage in je rapport opgenomen.

\documentclass{hogent-article}

% Voor een voorstel in het Engels: \documentclass[english]{hogent-article}
% Dat kan enkel na toestemming van je begeleider.

\addbibresource{voorstel.bib}

%%=============================================================================
%% Projectgegevens 
%%=============================================================================

%% ---------- Jij ------------------------------------------------------------
\newcommand{\studentnaam}{Sarah Turner Burkitt}
\newcommand{\studentemail}{sarah.turnerburkitt@student.hogent.be}

%% Wordt gebruikt als bestandsnaam van je PDF's, bijvoorbeeld
%% Familienaam-Voornaam-rapport.pdf. Geen spaties of leestekens gebruiken.
\newcommand{\bestandsnaam}{TurnerBurkitt-Sarah}

%% ---------- Je project -----------------------------------------------------
\newcommand{\projecttitel}{Gerichte configuratie van automatisch gegenereerde output}
\newcommand{\projectondertitel}{Uitbreiding van document- en consultatietemplates binnen Vesalius.ai}
%% Een of twee zinnen. Wordt gebruikt op de poster en in het voorstel.
\newcommand{\projectpitch}{%
    Dit graduaatsproject breidt het Vesalius-platform uit met gerichtere configuratie van automatisch gegenereerde output: 
    documenttemplates kunnen aan specifieke artsen worden gekoppeld en Scribe-consultatemplates kunnen worden afgestemd op patiënt- 
    of artsgerichte output.%
}

%% ---------- Waar je het doet ------------------------------------------------
\newcommand{\bedrijfsnaam}{Vesalius.ai}

%% Je mentor op de werkvloer (het bedrijf)
\newcommand{\mentornaam}{Bart D'Hoore}
\newcommand{\mentoremail}{Bart.D'Hoore@endare.be}

%% Je begeleider van de opleiding (HOGENT)
\newcommand{\begeleidernaam}{Luc Vervoort}
\newcommand{\begeleideremail}{luc.vervoort@hogent.be}

%% ---------- Je repository ---------------------------------------------------
\newcommand{\repolink}{https://github.com/STB95/Graduaatsproject_TurnerBurkittSarah_26-27}
%% op je poster: hij is een beoordeeld onderdeel van je oplevering.
\newcommand{\repolink}{https://les.lucvervoort.org/sarahturnerburkitt/graduaatsproject}

%% ---------- Vast voor dit opleidingsonderdeel -------------------------------
%% Deze twee hoef je normaal niet aan te passen.
\newcommand{\opleidingnaam}{Graduaat in het Programmeren}
\newcommand{\academiejaarwaarde}{2026-2027}

%% Zoekwoorden voor het voorstel en de poster (3 tot 5)
\newcommand{\projectkeywords}{Vesalius.ai, documenttemplates, Scribe, automatische output, Angular/Laravel, WYSIWYG e-maileditor}


%% Opleiding, vak en soort opdracht
\studyprogramme{\opleidingnaam}
\course{Graduaatsproject}
\assignmenttype{Projectvoorstel}

\academicyear{\academiejaarwaarde}

\title{\projecttitel}

\author{\studentnaam}
\email{\studentemail}

%% Je mentor op de werkvloer
\supervisor[Mentor]{\mentornaam{} (\bedrijfsnaam, \href{mailto:\mentoremail}{\mentoremail})}

%% De link naar je repository
\projectrepo{\repolink}

\keywords{\projectkeywords}

\begin{document}

\begin{abstract}
  % Eén samenhangende alinea van vijf tot acht zinnen, die je hele voorstel
  % samenvat: de context bij je bedrijf, het probleem, wat je gaat maken, met
  % welke technologie, en wat het bedrijf eraan heeft. Geen inleiding, geen
  % opsomming.
  Vervang deze tekst door de samenvatting van je voorstel.
\end{abstract}

%%=============================================================================
%% De inhoud van je projectvoorstel
%%=============================================================================
%%
%% Zeven rubrieken, samen hoogstens twee bladzijden. Houd het kort: per
%% rubriek volstaan één tot drie zinnen. Een voorstel is geen rapport - het
%% moet je begeleider in vijf minuten laten beslissen of dit een geschikt
%% graduaatsproject is.
%%
%% Stem je voorstel af met je mentor VOORDAT je het indient. Je mentor
%% bevestigt dat dit is wat het bedrijf bedoelt en dat de tien werkdagen
%% haalbaar zijn.

%---------- 1. Context en probleemstelling ------------------------------------

\section{Context en probleemstelling}%
\label{vst:context}

% Waar loop je stage, wat doet dat team, en welk concreet probleem of welke
% concrete behoefte is de aanleiding? Noem de mensen voor wie het een probleem
% is. Vermijd algemeenheden als "bedrijven hebben nood aan".

Vervang deze tekst.

%---------- 2. Doelstelling ---------------------------------------------------

\section{Doelstelling}%
\label{vst:doelstelling}

% Wat bestaat er op het einde van dit project dat er vandaag niet is? Formuleer
% het als een resultaat waarvan iemand kan vaststellen dat het er is, en
% vermeld waaraan je zal merken dat het geslaagd is.

Vervang deze tekst.

%---------- 3. Technologieën --------------------------------------------------

\section{Technologieën}%
\label{vst:technologieen}

% Met welke talen, frameworks, databanken en diensten ga je werken? Welke
% liggen vast door de omgeving van het bedrijf, en welke kies je zelf? Vermeld
% ook wat nieuw voor je is: technologische verbreding wordt gewaardeerd.
%
% Vermeld hier ook of je AI denkt in te zetten, en waarvoor. Dat mag en wordt
% zelfs verwacht; in je rapport verantwoord je later wat je precies gedaan hebt.
%
% Verwerkt je toepassing persoonsgegevens, of zit er AI in wat je bouwt? Zeg
% dat dan nu al in één zin: dan weet je begeleider meteen dat de AVG of de
% AI Act meespeelt, en kan dat mee bekeken worden bij de goedkeuring.

Vervang deze tekst.

%---------- 4. Afbakening -----------------------------------------------------

\section{Afbakening}%
\label{vst:afbakening}

% DE BELANGRIJKSTE RUBRIEK. Je stage en je graduaatsproject lopen bij hetzelfde
% bedrijf; dit is waar je ze uit elkaar houdt.
%
% Beschrijf:
%   - wat je dagelijkse stagewerk is;
%   - wat het afgebakende geheel is dat JIJ zelfstandig analyseert, ontwerpt
%     en oplevert - dat is je graduaatsproject;
%   - wat er uitdrukkelijk buiten valt.
%
% Toets voor jezelf: welke analyse en welk ontwerp maak jij zelf? Kan je daar
% niets over schrijven, dan heb je een taak en geen project.

Vervang deze tekst.

%---------- 5. Verwachte meerwaarde -------------------------------------------

\section{Verwachte meerwaarde}%
\label{vst:meerwaarde}

% Wat heeft het bedrijf hieraan? Wat blijft er van jouw werk over wanneer je
% stage afgelopen is? Denk aan tijdwinst, minder fouten, nieuwe mogelijkheden,
% lagere kosten, of een basis waarop verder gebouwd kan worden.

Vervang deze tekst.

%---------- 6. Planning en werkbelasting --------------------------------------

\section{Planning en werkbelasting}%
\label{vst:planning}

% Het graduaatsproject is ongeveer tien werkdagen, samen ongeveer 100 uren.
% Verdeel die over de fasen van je project en geef per fase aan wat je oplevert.
% Vermeld ook of je die dagen gefocust opneemt of gespreid over de stageperiode,
% en dat je dit met je mentor afgesproken hebt.

\begin{table}[htbp]
  \centering
  \begin{tabular}{llr}
    \toprule
    \textbf{Fase} & \textbf{Wat lever je op} & \textbf{Uren} \\
    \midrule
    Analyse en vereisten      & overzicht van de vereisten   & 15  \\
    Technologiekeuze          & verantwoorde keuze           & 10  \\
    Ontwerp                   & architectuur en gegevensmodel& 15  \\
    Realisatie                & werkende oplossing           & 40  \\
    Testen en validatie       & testverslag                  & 10  \\
    Rapport, poster en demo   & de op te leveren documenten  & 10  \\
    \midrule
    \textbf{Totaal}           &                              & \textbf{100} \\
    \bottomrule
  \end{tabular}
  \caption[Planning en werkbelasting]{\label{tab:vst-planning}%
    Vervang deze verdeling door je eigen planning.}
\end{table}

%---------- 7. Risico's -------------------------------------------------------

\section{Risico's}%
\label{vst:risicos}

% Wat kan dit project doen mislukken, en wat doe je dan? Denk aan: toegang tot
% gegevens of systemen die je nodig hebt, afhankelijkheid van een collega of
% een derde partij, technologie die je nog moet leren, of drukte op de
% werkvloer waardoor je dagen niet vrijkomen.
%
% Twee tot drie risico's met telkens één zin over wat je terugvalplan is.

Vervang deze tekst.


\printbibliography[heading=bibintoc]

\end{document}
